\chapter{Who We Target First}

We are not targeting ``everyone using AI.''
We are targeting teams that already feel the pain of unexplained AI decisions.

The key filter is not company size.
The key filter is urgency.

\section{Ideal customer profile}

A good first customer has all three of the following:

\begin{enumerate}
  \item AI is already live in a real workflow.
  \item The workflow affects money, customers, or trust.
  \item Something has already gone wrong, or the team can see that it will.
\end{enumerate}

\begin{conceptbox}{Three-question qualifier}
Ask these three questions before investing time in a prospect:

\smallskip
\begin{enumerate}
  \item Is your AI already taking actions, not just producing text?
  \item When it makes the wrong call, can you explain why?
  \item Have you had an incident yet, or do you expect one?
\end{enumerate}

If the answer is yes, yes (they cannot), and yes --- that is our customer.
\end{conceptbox}

\section{Segment 1 --- AI support startups and support automation teams}

This is the clearest early wedge.

The workflow is easy to understand, the consequences are visible, and a bad AI decision shows up fast.

Real pain examples:

\begin{itemize}
  \item a refund is issued incorrectly and the team cannot explain why,
  \item a discount was applied too aggressively without a clear decision trace,
  \item a customer escalation cannot be reconstructed from logs,
  \item a case was closed or routed the wrong way and the team has no trail.
\end{itemize}

\textbf{Who pays:} typically the VP of Support, Head of CX, or a technically strong CTO.

\textbf{Why they pay:} they are losing time in incident review, losing customer trust, and in regulated industries they need defensible output.

\textbf{Sales cycle:} short. If the pain is real and the workflow is live, a pilot conversation can open within days.

\section{Segment 2 --- Fintech ops teams}

These buyers care because AI decisions in their world touch money, risk, and regulatory surface.

Real pain examples:

\begin{itemize}
  \item transaction review where AI flags or passes decisions without a clear trace,
  \item fraud ops where decisions need to be defensible to internal or external auditors,
  \item risk scoring workflows where model behavior has to be explainable,
  \item exception handling where a human needs to understand why the AI routed a case.
\end{itemize}

\textbf{Who pays:} Head of Ops, Head of Risk, or a fintech founder in a startup context.

\textbf{Why they pay:} compliance and audit risk. The cost of one unexplained AI decision in a regulated fintech is higher than the cost of the tool.

\textbf{Sales cycle:} medium. There may be a procurement or security review, but the pain is real enough to move if we frame it correctly.

\section{Segment 3 --- Internal automation teams}

These are teams using AI inside the company for approvals, case routing, exception handling, and operations.

This segment is good only when the workflow already matters and has already caused friction.

Real pain examples:

\begin{itemize}
  \item AI-powered procurement routing with no trace when a case goes wrong,
  \item claims or exception handling where humans are pulled in to explain AI output,
  \item internal approvals driven by AI with no audit log available.
\end{itemize}

\textbf{Who pays:} the department head or an internal automation engineering team.

\textbf{Why they pay:} internal trust. Other teams will not rely on an automation they cannot understand or verify.

\section{Buyer mindset}

These buyers are \textbf{not} shopping for AI infrastructure.
That language is too abstract.

They are shopping for:

\begin{itemize}
  \item control,
  \item visibility,
  \item safety,
  \item explanation,
  \item proof.
\end{itemize}

\begin{insightbox}{How to position to the buyer}
Do not say: ``Connector is a governed agentic runtime.''

Say: ``When your AI makes a decision, Connector shows what happened, why it happened, and gives you evidence you can stand behind.''
\end{insightbox}

That framing is concrete, relatable, and close to a buying decision.
